\section{So sánh và lựa chọn phong cách kiến trúc phù hợp (Compare and choose suitable architecture styles)}
Việc chọn architecture style cho IRMS được xem như một bài toán trade-off giữa các \textit{architecture characteristics} (non-functional requirements). IRMS không chỉ có nhiều chức năng (menu, order, kitchen workflow, table/reservation, billing, inventory, reporting) mà còn yêu cầu vận hành trơn chu giữa front-of-house và kitchen, tức là phải giảm delay/error và đồng bộ trạng thái theo thời gian gần realtime. Vì vậy, các tiêu chí như performance, availability/reliability, scalability, maintainability/extensibility và security được ưu tiên khi so sánh các style.

\subsection{Modular Monolith.}
Theo ma trận đánh giá, Modular Monolith có điểm mạnh lớn nhất ở simplicity và cost thấp. Với bối cảnh đó, cách tiếp cận này giúp nhóm build nhanh, debug dễ và hạn chế overhead vận hành (deploy 1 lần thay vì nhiều service). Tuy nhiên, worksheet cũng chỉ ra monolith bị hạn chế ở scalability, elasticity và fault-tolerance. Đối với IRMS, giờ cao điểm có thể tạo bottleneck ở các luồng như ``place order $\rightarrow$ send to kitchen $\rightarrow$ update status $\rightarrow$ checkout''. Nếu toàn bộ hệ thống chạy chung một khối, khi một module (ví dụ kitchen queue hoặc payment) gặp sự cố hoặc bị quá tải, hiệu ứng lan truyền có thể làm giảm availability của toàn hệ thống. Đây là trade-off quan trọng cần cân nhắc.

\subsection{Microservices.}
Microservices được đánh giá cao nhất về maintainability, testability, deployability, scalability, fault-tolerance và evolvability. Điều này phù hợp với IRMS vì các module có vòng đời thay đổi khác nhau: menu/promotion thường xuyên update, kitchen workflow cần tối ưu liên tục, billing/payment có tính nhạy cảm và hay phụ thuộc third-party, còn reporting có thể xử lý theo batch. Khi tách thành các component/service độc lập, việc scale và deploy từng phần sẽ chủ động hơn (ví dụ scale riêng kitchen hoặc ordering trong giờ cao điểm), đồng thời việc refactor hoặc nâng cấp cũng ít ảnh hưởng dây chuyền. Trade-off ở đây là chi phí triển khai và vận hành tăng (nhiều service, nhiều API call giữa service, cần cơ chế quan sát/monitoring tốt), tức là overhead sẽ lớn hơn so với monolith.

\subsection{Event-driven.}
Event-driven architecture được worksheet đánh giá nổi bật ở responsiveness, fault-tolerance và evolvability. Với IRMS, nhiều tình huống tự nhiên đã mang tính sự kiện: OrderPlaced/OrderUpdated, ticket gần deadline, món Done/Ready-to-serve, payment success/fail, low-stock alert, audit log... Nếu các tương tác này được handle bằng event (publish/subscribe) thay vì gọi đồng bộ trực tiếp, các component sẽ giảm coupling, giảm rủi ro chờ đợi khi một service chậm, và dễ mở rộng thêm hành vi (ví dụ thêm module analytics lắng nghe sự kiện để tính SLA) mà không cần sửa code các module hiện có. Trade-off là hệ thống phức tạp hơn ở phần consistency và xử lý lỗi: cần thiết kế idempotency (tránh duplicate event), cơ chế retry/replay khi hệ thống bếp offline, và chấp nhận một mức eventual consistency cho một số trạng thái.

\subsection{Kết luận lựa chọn.}
Từ các phân tích trên, báo cáo này định hướng lựa chọn kiến trúc theo hướng kết hợp: dùng phân rã theo các ``architecture quantum'' có thể deploy độc lập (phù hợp với microservices) cho các domain chính như Ordering, Kitchen, Billing, Inventory, Reporting; đồng thời các luồng liên module cần realtime và ít coupling (đặc biệt là Ordering $\leftrightarrow$ Kitchen theo UC6--UC9) sẽ được thiết kế theo event-driven để tăng responsiveness và fault-tolerance. Khi thiết kế các component, anti-pattern ``entity trap'' cần được tránh (không tạo các lớp/manager khổng lồ theo từng entity), thay vào đó ranh giới module/component nên bám theo chức năng và workflow thực tế (ví dụ: kitchen ticket workflow + state machine) nhằm giữ cohesion cao và hỗ trợ maintain/extensibility theo yêu cầu của hệ thống.